\addtocontents{toc}{\protect\setcounter{tocdepth}{1}}
\chapter[Pruebas]{Pruebas}

\emph{Adaptación basada en IEEE Std 829-1998, Software Test Documentation
\citep{ieee_829_1998}.}

\section[Elementos]{Elementos de prueba}

Las pruebas verifican los siguientes elementos del sistema.

\begin{itemize}
    \item \emph{Pantallas}: acceso y recuperación de cuenta, panel principal,
    registro clínico, historial, chat, preguntas frecuentes, biblioteca, detalle
    de material, reproductor de video y recordatorios.
    \item \emph{Motor de reglas clínicas}: clasificación de cada registro en los
    niveles normal, alerta y crítico.
    \item \emph{Servicio de cifrado}: cifrado y descifrado de los datos
    sensibles.
    \item \emph{API del servidor}: autenticación, aislamiento entre usuarios,
    rechazo de campos clínicos en claro, escritura idempotente, paginación,
    verificación de correos y plantillas de correo.
    \item \emph{Exportadores}: generación de los archivos PDF y Excel.
    \item \emph{Servicios de caché y notificaciones}: disponibilidad del
    contenido sin conexión y programación de recordatorios.
\end{itemize}

\section[Especificación]{Especificación de las pruebas}

La Tabla~\ref{tab:especificacion-pruebas} resume la estrategia para cada nivel
de prueba.

\begin{table}[H]
\centering
\small
\tablazebra
\caption{Especificación de las pruebas por nivel.}
\label{tab:especificacion-pruebas}
\begin{tabularx}{\textwidth}{P{0.62}P{1.19}P{1.19}}
\toprule
\filaencabezado \encab{Aspecto} & \encab{Pruebas de unidad} & \encab{Pruebas de sistema} \\
\midrule
Características & Funcionalidad y seguridad de servicios y reglas. &
Funcionalidad, interfaz y navegación de cada pantalla. \\
Objetivo & Detectar errores de lógica en el motor de reglas, el cifrado y el
acceso a datos. & Verificar que cada pantalla cumple los criterios de aceptación
de su historia de usuario. \\
Enfoque & Caja blanca. & Caja negra. \\
Técnicas & Valores límite en los umbrales de signos vitales y en la escala ESAS,
y particiones de equivalencia por nivel de alerta. & Casos derivados de los
criterios de aceptación y verificación de pantallas en distintos tamaños. \\
Actividades & Preparar datos simulados, ejecutar flutter test y npm test en el
backend, y corregir las fallas. & Montar cada pantalla con proveedores simulados,
interactuar con ella y comparar el resultado esperado; recorrer la API real con
el script de humo y la prueba manual en un teléfono. \\
Criterio de cumplimiento & El 100~\% de los casos aprobados. & El 100~\% de
los casos aprobados, sin desbordes de pantalla. \\
\bottomrule
\end{tabularx}
\end{table}

Las pruebas de la aplicación se ejecutan de forma automatizada con flutter
test, con dobles de prueba para el cliente de la API, de modo que no dependen
de la red. Las pruebas del servidor se ejecutan con Jest y supertest contra una
base PostgreSQL de pruebas real, levantada con Docker Compose, por lo que
verifican las consultas y las restricciones de la base y no un simulacro de
ellas. Además, el script tool/humo\_api.dart recorre la API real de punta a
punta: registro, clave de datos e inicio de sesión.

\section[Responsables]{Responsables de las pruebas}

Las dos estudiantes desarrolladoras diseñan, ejecutan y mantienen la suite, y
corrigen los defectos que detecta. La contraparte de
Enfermería revisa en cada entrega la pertinencia de los textos, los contenidos y
los criterios de alerta.

\section[Calendario]{Calendario de pruebas}

Las pruebas acompañan a cada sprint. Toda historia de usuario se integra con sus
pruebas, y al cierre del sprint se ejecuta la suite completa antes de generar el
APK de revisión.

\begin{table}[H]
\centering
\small
\tablazebra
\caption{Calendario de pruebas por sprint.}
\label{tab:calendario-pruebas}
\begin{tabularx}{\textwidth}{P{0.50}P{0.80}P{1.70}}
\toprule
\filaencabezado \encab{Sprint} & \encab{Periodo} & \encab{Pruebas} \\
\midrule
Sprint 1 & 26 ago -- 15 sep & Acceso, cifrado, registro clínico, motor de
reglas e historial. \\
Sprint 2 & 16 sep -- 20 oct & Panel, exportación, biblioteca, contenido sin
conexión, video, chat, preguntas frecuentes, conversaciones, recordatorios. \\
Sprint 3 & 21 oct -- 13 nov & Datos por paciente, aislamiento entre usuarios,
navegación y sincronización sin conexión. Ejecución completa de la suite. \\
\bottomrule
\end{tabularx}
\end{table}

\section[Detalle de las Pruebas]{Detalle de las pruebas}

La suite de la aplicación suma 294 casos de prueba en 36 archivos, y la del
servidor, 54 casos en 4 archivos, para un total de 348 casos automatizados
(cifras de la ejecución más reciente). Las carpetas de prueba de la aplicación
replican la estructura del código: cada funcionalidad tiene sus pruebas de
dominio, de datos y de presentación. La Tabla~\ref{tab:inventario-tests}
relaciona cada área con los archivos que la verifican y los requerimientos que
cubre; los archivos se nombran sin la ruta ni el sufijo \emph{\_test}.

\begin{table}[H]
\centering
\small
\tablazebra
\caption{Distribución de la suite automatizada.}
\label{tab:inventario-tests}
\begin{tabularx}{\textwidth}{P{0.75}P{0.30}P{1.35}P{0.60}}
\toprule
\filaencabezado \encab{Área} & \encab{Archivos} & \encab{Elemento verificado} & \encab{RF} \\
\midrule
App (navegación y tema) & 5 & navegacion\_principal, fundido\_ramas y paleta del tema con su contraste & RF-18 \\
Autenticación (dominio) & 1 & validaciones\_registro & RF-02 \\
Biblioteca & 10 & Categorías, secciones, temas y modelo del material, caché de contenido, reproductor, visor de PDF y tarjeta & RF-24 a RF-27 \\
Chat & 2 & Modelo de conversación y burbujas del chat & RF-19, RF-20 \\
Configuración & 2 & Escala de texto y modo de tema & --- \\
Historial & 2 & Exportadores y vista previa de Excel & RF-16, RF-17 \\
Pacientes & 1 & proveedores\_pacientes (paciente activo) & RF-08 \\
Panel y preguntas frecuentes & 2 & pantalla\_panel\_principal y pantalla\_preguntas\_frecuentes & RF-18, RF-21 \\
Registro clínico (dominio) & 2 & motor\_reglas\_clinicas y registro\_clinico & RF-11, RF-13 \\
Compartido & 6 & Widgets reutilizables, entre ellos el banner de conexión & RF-28 \\
Núcleo & 3 & servicio\_conectividad, calendario\_avisos y silencio\_avisos & RF-23, RF-28 \\
\midrule
\filatotal \textbf{Aplicación (flutter test)} & \textbf{36} & \textbf{294 casos} & \\
\midrule
Servidor: auth & 1 & Registro, inicio de sesión, rutas protegidas, clave de datos y recuperación de contraseña & RF-02 a RF-05 \\
Servidor: datos & 1 & Aislamiento entre usuarios, rechazo de campos en claro, upsert idempotente, paginación, cascadas, conversaciones y favoritos & RF-07, RF-10, RF-11, RF-14, RF-20, RF-22, RF-28 \\
Servidor: correos & 1 & Cambio de correo, correo de respaldo y verificación del registro con enlaces de un solo uso & RF-04, RF-06 \\
Servidor: plantillas & 1 & Plantillas de correo y páginas web de restablecimiento y confirmación & RF-04, RF-06 \\
\midrule
\filatotal \textbf{Servidor (npm test)} & \textbf{4} & \textbf{54 casos} & \\
\bottomrule
\end{tabularx}
\end{table}

Las pruebas del servidor son las que sostienen las garantías de seguridad del
sistema. Verifican que un usuario no ve ni modifica datos de otro, que el
servidor rechaza con error 400 los campos clínicos enviados en claro, que una
escritura repetida no duplica registros, que la paginación por fecha y por rango
entrega los registros esperados y que los enlaces de verificación vencen, son de
un solo uso y no revelan si un correo existe.

Como evidencia complementaria, el sistema se recorrió de punta a punta en un
teléfono Android conectado a la API: registro de la cuenta, inicio de sesión,
alta de un paciente y envío y confirmación de los correos de verificación. Esta
prueba manual no genera un registro de resultados por caso.

El motor de reglas clínicas concentra las pruebas más detalladas de la aplicación, porque de él
depende la clasificación del nivel de alerta. Sus casos recorren los
umbrales de cada signo vital y de la escala ESAS con valores límite, por ejemplo
una temperatura de 39,5~°C, que todavía se clasifica como alerta, frente a una
de 39,6~°C, que ya es crítica. El Anexo de especificación de pruebas detalla los
casos por unidad.

\section[Conclusiones]{Conclusiones de prueba}

Las suites de la aplicación y del servidor cubren la mayoría de los
requerimientos funcionales y permiten ejecutar la verificación completa en pocos minutos
antes de cada entrega. El énfasis en el motor de reglas y en el cifrado responde
a que ambos concentran el riesgo del sistema, uno por su efecto en la
orientación al cuidador y el otro por la protección de datos de salud.

La suite incluye también los requerimientos del sprint~3, implementados por
adelantado: la separación de datos por paciente, el aislamiento entre usuarios y
la sincronización sin conexión. Queda pendiente validar en un dispositivo real
la caída y la recuperación de la red, que las pruebas automáticas simulan.

\addtocontents{toc}{\protect\setcounter{tocdepth}{2}}
